\documentclass{article}

\usepackage{../../Style}

\title{Differential Geometry}
\author{Shing Tak Lam}
\date{Michaelmas 2023\footnote{Based on lectures by Alexei Kovalev. Last updated \today.}}

\makeindex

\begin{document}

\maketitle

\tableofcontents

\lec

% Four examples classes, given by A.G.Kovalev (Oct 26 etc.). Note last examples class is Thursday Week 9.

\section{Manifolds}

The concept of a manifold is intended to generalise curves and surfaces in \(\R^3\).

% Motivation
% Definition of a topological space
% Definition of Huasdorff and second countable

\begin{definition}
    [smooth structure]
    A \emph{smooth structure}\index{smooth structure} on a topological space \(M\) is a collection of \emph{charts}\index{charts} (an \emph{atlas}\index{atlas})
    \[\varphi_\alpha : U_\alpha \to V_\alpha\]
    where \(U_\alpha \sse M\), \(V_\alpha \sse \R^d\) are open, \(\varphi_\alpha\) is a homeomorphism, such that
    \begin{enumerate}
        \item \(M = \bigcup_\alpha U_\alpha\). That is, the charts cover \(M\).
        \item Given any charts \(\varphi_\alpha, \varphi_\beta\), the map
        \[\varphi_\beta \circ \varphi_\alpha^{-1} : \varphi_\alpha(U_\alpha \cap U_\beta) \to \varphi_\beta(U_\alpha \cap U_\beta)\]
        is a smooth map between open subsets of \(\R^d\).
        \item If \(\varphi\) is \emph{compatible} (as in the previous point) with every chart in the atlas, then \(\varphi\) is a chart in the atlas.
    \end{enumerate}
\end{definition}

In practice (i.e. using Zorn's lemma), given any collection of charts satisfying the first two conditions, we can find a maximal atlas containing it.

\begin{remark}
    \begin{enumerate}
        \item The second condition implies that \(\varphi_\beta\circ \varphi_\alpha^{-1}\) is a diffeomorphism.
        \item \(d\) is fixed, assuming that \(M\) is connected.
    \end{enumerate}
\end{remark}

\begin{definition}
    [manifold] A \emph{(smooth) manifold}\index{manifold} \(M\) is a second countable Huasdorff space topological space, equipped with a smooth structure.

    \(d = \dim(M)\) is the \emph{dimension}\index{dimension} of \(M\).
\end{definition}

In practice, we may induce a topology on \(M\) from a smooth structure (without the mention of open or homeomorphism), then \(D \sse M\) is open if for any chart \(\varphi_\alpha\), \(\varphi_\alpha(D \cap U_\alpha)\) is an open subseteq of \(\R^d\).

\begin{remark}
    \begin{enumerate}
        \item We can also define a \(C^k\) manifold by replacing smooth with \(C^k\) as above. If \(k = 0\), then we have a \emph{topological manifold}\index{topological manifold}.
        \item On the other hand, we can replace \(\R^d\) with \(\C^n\), and obtain a \emph{complex manifold}\index{complex manifold}.
    \end{enumerate}
\end{remark}

\(\R^n\) is (trivially) a manifold.

\begin{example}
    [unit sphere]
    Define
    \[S^n = \left\{x = (x_0, \dots, x_n) \in \R^n \mid \norm{x} = 1\right\} \sse \R^{n+1}\]
    For the charts, we can consider the stereographic projections
    \[\varphi(x) = \frac{1}{1-x_0}(x_1, \dots, x_n) \qquad \text{for }x_0 \ne 1\]
    and
    \[\psi(x) = \frac{1}{1+x_0}(x_1, \dots, x_n) \qquad \text{for }x_0 \ne -1\]
    In particular,
    \[\psi \circ \varphi^{-1}(u) = \frac{u}{\norm{u}^2}\]
\end{example}

\begin{example}
    If \(M, N\) are manifolds of dimensions \(d,e\) respectively, then \(M \times N\) is a manifold of dimension \(d + e\).
\end{example}

As a corollary, \(T^n = (S^1)^n\) is a manifold.

\begin{example}
    If \(M\) is a manifold, \(X \subseteq M\) is open, then \(X\) is a manifold.
\end{example}

\lec

% Example sheet 1 out, will mark questions 1 and 2 (via moodle)
% Example class 1 - Thursday 26 Oct 3:30pm - 5pm MR13

A convention, we will will assume for any manifold \(M\), we will assume all components are of the same dimension. In particular, \(\dim(M)\) is well defined.

\begin{notation}
    We will write \(M^d\) to say \(M\) is a manifold of dimension \(d\).
\end{notation}

\begin{example}
    [real projective space] The real projective space\index{real projective space@real projective space \(\R\P^n\)} is
    \[\R\P^n = \left\{\text{straight lines in }\R^{n+1}\text{ through }0\right\}\]
    Points in \(\R\P^n\) can be written using homogeneous coordinates \((x_0 : \dots : x_n)\), where the \(x_i\) are not all zero, and
    \[(x_0 : \dots : x_n) = (\lambda x_0 : \dots : \lambda x_n)\]
    for all \(\lambda \ne 0\). We will induce a topology from a smooth structure. The charts are given by \((\varphi_i, U_i)\) where
    \[U_i = \left\{x_i \ne 0\right\}\]
    and we define
    \[\varphi_i(x_0:\dots:x_n) = \left(\frac{x_0}{x_i}, \dots, \frac{x_{i-1}}{x_i},\frac{x_{i+1}}{x_i}, \dots, \frac{x_n}{x_i}\right) \in \R^n\]
    For \(i < j\), 
    \[\varphi_j \circ \varphi_i^{-1}(y_0, \dots, y_n) = \left(\frac{y_1}{y_j}, \dots, \frac{1}{y_j}, \dots, \frac{y_n}{y_j}\right)\]
    where the \(1\) is in the \(i\)-th position, and we omit the \(j\)-th position. Thus, \(\R\P^n\) is an \(n\)-dimensional manifold. Similarly, we can define the complex projective space \(\C\P^n\)\index{complex projective space@complex projective space \(\C\P^n\)}, which is a \(2n\)-manifold (and a \(n\)-dimensional complex manifold).
\end{example}

\begin{example}
    [Grassmannian]
    \[\Gr(k, n) = \left\{k\text{-dimensional subspaces of }\R^n\right\}\]\index{Grassmannian@Grassmannian \(\Gr(k, n)\)}
    Then \(\Gr(k, n)\) is a manifold of dimension \(k(n - k)\). Note we can also define the Grassmannian over \(\C\).
\end{example}

\begin{proof}
    [Sketch proof] An example of a coordinate neighbourhood is
    \[U = \left\{k\text{-dimensional subspaces obtainable as the span of rows of }k \times n\text{ matrices of the form} \begin{pmatrix}
        I_k & *
    \end{pmatrix}\right\}\]
    where \(I_k\) is the \(k \times k\) identity matrix. In this case, the \(k \times (n-k)\) block \(*\) defines local coordinates. Choosing different columns, we have \({n \choose k}\) neighbourhoods, \(U_{i_1, \dots, i_k}\), where \(i_1 < i_2 < \dots < i_k\).
\end{proof}

In particular, note that \(\R\P^n = \Gr(1, n+1)\).

\begin{example}
    [Non-example] Consider an equivalence relation on \(\R^2\) given by
    \[(x, y) \sim \left(\lambda x, \frac{1}{\lambda}y\right)\]
    for any \(\lambda \ne 0\). Let \(X = \R^2/\sim\) for the quotient, with the quotient topology. We can consider
    \[\left\{(x, y) \mid xy = c\right\}\]
    If \(c \ne 0\), then this is one equivalence class. If \(c = 0\), this splits into three, which are
    \begin{itemize}
        \item \(\{(0, 0)\}\),
        \item \(\{(x, 0) \mid x \ne 0\}\),
        \item \(\{(0, y) \mid y \ne 0\}\)
    \end{itemize}
    We can write
    \[X = (-\infty, 0) \cup (0, \infty) \cup \left\{0', 0'', 0'''\right\}\]
    We have three charts \(\varphi_i\) corresponding to each choice of zero. For example,
    \[U_1 = (-\infty, 0) \cup \{0'\} \cup (0, \infty)\]
    with the ``natural'' map \(\varphi_1\). On the other hand, the induced topology is not Hausdorff. Note the induced topology is \emph{not} the quotient topology.
\end{example}

Refer to Examples Sheet 1 Question 12 for an example of a non-second countable (non-)example.

\begin{definition}
    [smooth map] Let \(M, N\) be manifolds, a continuous map \(f : M \to N\) is \emph{smooth}\index{smooth map} if for any \(p \in M\), we can find charts \((\varphi, U)\) near \(p\), \((\psi, V)\) near \(f(p)\), and the composition
    \[\psi \circ f \circ \varphi^{-1}\]
    is smooth. as a map on
    \[\varphi(U \cap f^{-1}(V))\]
\end{definition}

\begin{definition}
    [diffeomorphism] A smooth map \(f : M \to N\) is a \emph{diffeomorphism}\index{diffeomorphism} if \(f\) is a bijection, and \(f^{-1}\) is smooth. In this case, we say that \(M, N\) are \emph{diffeomorphic}\index{diffeomorphic}.
\end{definition}

\begin{proposition}
    \begin{enumerate}
        \item If \(U \sse \R^n\), then \(f : U \to \R^m\) is smooth in the above definition if and only if it is smooth in the sense of multi-variate calculus.
        \item Every chart \(\varphi : U \to \varphi(U) \sse \R^d\) is a smooth map of manifolds. Moreover, it is a diffeomorphism.
        \item If \(f : M \to N\), \(g : N \to P\) are smooth, then \(g \circ f : M \to P\) is smooth.
    \end{enumerate}
\end{proposition}

\begin{proof}
    Obvious, and so omitted.
\end{proof}

\subsection{Matrix Lie groups}

Consider the group \(\GL(n, \R) \sse \R^{n^2}\). In fact, this is an open subset, and so it is a manifold. Moreover, matrix multiplication and inversion are smooth.

\begin{definition}
    [Lie group] A group \(G\) is called a \emph{Lie group}\index{Lie group} if
    \begin{itemize}
        \item \(G\) is a manifold,
        \item the multiplication map \(G \times G \to G\) is smooth,
        \item inversion \(G \to G\) is smooth.
    \end{itemize}
\end{definition}

\begin{remark}
    Suffices to show that the map
    \[G \times G \ni (\sigma, \tau) \mapsto \sigma\tau^{-1} \in G\]
    is smooth.
\end{remark}

In particular, \(\GL(n, \R\) is a Lie group. Similarly, \(\GL(n, \C)\) is also a Lie group. More generally, for \(A = (a_{ij}) \in \Mat_n(\C)\), we define the norm \(\norm{A} = n\max_{i, j}\abs{i, j}\). (Note in finite dimensional, all norms are equivalent). In particular,
\[\norm{AB} \le \norm{A}\norm{B}\]
We define the \emph{exponential of \(A\)}\index{exponential map!of matrices} as
\[\exp(A) = \sum_k \frac{A^k}{k!}\]
This series converges for all \(A\), by the same proof as for \(\exp : \C \to \C\) (or more generally, any Banach algebra). Moreover, it converges uniformly on balls
\[\left\{\norm{A} \le R\right\}\]
by the Weierstrass \(M\)-test. In particular, \(\exp : \Mat_n(\C) \to \Mat_n(\C)\) is continuous.

\lec

First of all, note that
\[f(A) = A^n\]
is a differentiable function, with
\[\norm{(\dd f)_A} \le n\abs{A}^{n-1}\]
Therefore, the series which we get by differentiating term by term, converges uniformly on any bounded set, again by the Weierstrass \(M\)-test.

Therefore, \(\exp(A)\) is (at least) \(C^1\) on \(\Mat_n(\C)\). We can similarly consider the higher derivatives, and we find that \(\exp(A)\) is \(C^\infty\).

\begin{remark}
    Some basic properties of the matrix exponential:

    \begin{enumerate}
        \item \(\exp(A^\T) = \exp(A)^\T\),
        \item \(\exp(A^*) = \exp(A^*)\),
        \item \(\exp(CAC^{-1}) = C\exp(A)C^{-1}\),
    \end{enumerate}

    In general,
    \[\exp(A + B) = \exp(A)\exp(B)\]
    does not hold. However, it does hold if \(AB = BA\). In particular,

    \begin{enumerate}
        \setcounter{enumi}{3}
        \item \(\exp(A)\exp(-A) = \exp(-A)\exp(A) = I\)
    \end{enumerate}
\end{remark}

We can also define the \emph{matrix logarithm}\index{logarithm map!for matrices} via the series
\[\log(I + A) = \sum_m(-1)^{m+1}\frac{A^m}{m}\]
We can show that this converges absolutely and uniformly on \(\{\abs{A} < \epsilon\}\) for any \(0 < \epsilon < 1\). The same result holds for the term by term differentiation, and again we get that \(\log(I + A)\) is \(C^\infty\) on \(\{\abs{A} < 1\}\)

Suppose \(\log(A)\) is defined, that is, \(\abs{A - I} < 1\), then
\[\exp(\log(A)) = A\]
since the same proof as in \(\C\) works for any Banach algebra.

On the other hand, if we would like
\[\log(\exp(A)) = A\]
we clearly need \(\abs{\exp(A) - I} < 1\). Suppose
\[A_\theta = \begin{pmatrix}
    0 & -\theta \\
    \theta & 0
\end{pmatrix}\]
for some \(\theta \in \R\). In particular,
\[A_\theta^2 = -\theta^2I\]
and we find that
\[\exp(A_\theta) = \begin{pmatrix}
    \cos(\theta) & -\sin(\theta) \\
    \sin(\theta) & \cos(\theta)
\end{pmatrix}\]
Setting \(\theta = 2\pi\), then we find that \(\exp(A_{2\pi}) = I\), but
\[\log(\exp(A_{2\pi})) = \log(I) = 0 \ne A_{2\pi}\]
Hence \(\abs{\exp(A) - I} < 1\) is a necessary, but not sufficient. On the other hand, \(\abs{A} < \log(2)\) is sufficient, since in this case,
\[\abs{e^\abs{A} - 1} < 1\]
which is required for absolute convergence.

\begin{example}
    [orthogonal group]
    Consider
    \[O(n) = \left\{A \in \GL(n, \R) \mid AA^\T = A^\T A = I\right\}\]
    We will show that \(O(n)\)\index{orthogonal group, \(O(n)\)} is a Lie group. Let \(A\) be an orthogonal matrix, \(\abs{A - I} < \epsilon < 1\), then \(B = \log(A)\) is defined. Moreover, \(\exp(B) = A\). By continuity, we can take \(\epsilon\) small enough such that \(\abs{B} < \log(2)\). Now
    \[\exp(B)\exp(B^\T) = AA^\T = \id\]
    and so
    \[\exp(B) = A = (A^\T)^{-1} = \exp(B^\T)^{-1} = \exp(-B^\T)\]
    Taking \(\log\), since \(\abs{B} = \abs{B^\T} < \log(2)\), we see that \(B = -B^\T\).

    Conversely, if \(B\) is skew-symmetric, with \(\abs{B} < \log(2)\), then
    \[\exp(B)\exp(B^\T) = I\]

    Define
    \[V_0 = \left\{B \in \Mat_n(\R) \mid  \abs{B} < \log(2), B^\T + B = 0\right\}\]
    and define
    \[U = \exp(V_0)\]
    which is an open neighbourhood of \(I \in O(n)\).

    \begin{proposition}
        \(O(n)\) has a \(C^\infty\) structure, making it into a manifold and a Lie group, with dimension \(n(n-1)/2\).
    \end{proposition}

    \begin{proof}
        We will use the subspace topology induced from \(O(n) \sse \R^{n^2}\). Define
        \begin{align*}
            h : U &\to V_0 \\
            h(A) &= \log(A)
        \end{align*}
        This is a homeomorphism onto its image, which is an open neighbourhood of \(0 \in \so(n) \cong \R^{n(n-1)/2}\).

        Given \(C \in O(n)\), define
        \[U_C = \left\{CA \mid A \in U\right\}\]
        and
        \begin{align*}
            h_C : U_C &\to V_0 \\
            h_C(A) &= \log(C^{-1}A)
        \end{align*}
        Again, this is a homeomorphism onto its image. Since \(C \in U_C\), we have an open cover. Moreover,
        \[h_{C_2} \circ h_{C_1}^{-1}(B) = h_{C_2}(C_1\exp(B)) = \log(C_2^{-1}C_2\exp(B))\]
        which is a smooth map between open neighbourhoods of \(\so(n)\). Therefore, \(O(n)\) is a smooth manifold. Similarly, for consider the map
        \[F(A_1, A_2) = A_1A_2^{-1}\]
        In local coordinates,
        \begin{align*}
            F_{\text{loc}}(B_1, B_2) &= h_{C_1C_2^{-1}}(F(h_{C_1}^{-1}(B_1), h_{C_2}^{-1}(B_2))) \\
            &= \log((C_1C_2^{-1})^{-1}C_1\exp(B_1)(C_2\exp(B_2))^{-1}) \\
            &= \log(C_2\exp(B_1)\exp(-B_2)C_2^{-1})
        \end{align*}
        and we can see that this is a smooth map as a function of \(B_1, B_2\).
    \end{proof}
\end{example}

This method works for more groups, and is sometimes called the Cayley construction (for matrix Lie groups).

\section{Tangent spaces to manifolds}

Consider a local curve in a ball in \(\R^n\), defined by a smooth parametrisation \(x(t) = (x_1(t), \dots, x_n(t))\). Suppose we have that \(x(0) = p \in \R^n\). In this case, its tangent vector at \(p\) is
\[\dot x(0)\in \R^n\]
If \(y(x)\) is a smooth change of coordinates near \(p\), then by the chain rule, we have
\[\dv{t}\bigg\vert_{t=0}y(x(t)) = J(y)_p \dot x(0)\]
where \(J(y)\) is the Jacobian matrix. More explicitly,
\[\dot y_i(0) = \sum_j \pdv{y_j}{x_j}(p)\dot x_j(0)\]

\begin{definition}
    [tangent vector]
    A \emph{tangent vector}\index{tangent vector} \(a\) to a manifold \(M^n\) at a point \(p \in M\), is an assignment to each coordinate chart \((U, \varphi)\) with \(p \in U\), an \(n\)-tuple of coordinates \((a_1, \dots, a_n) \in \R^n\), such that if \((U', \varphi')\) is another chart near \(p\), \((x_i), (x_i')\) the respective local coordinates,
    \[a_i' = \sum_j \pdv{x_i'}{x_j}(p)a_j\] 
\end{definition}

This definition is called the \emph{tensorial definition}, where we focus on the transformation law. There are several other definitions, for example, derivations, equivalence of curves to first order and so on.

\begin{definition}
    [tangent space]
    The \emph{tangent space}\index{tangent space} to \(M\) at a point \(p\), denoted \(\T_pM\) is the space of all tangent vectors to \(M\) at \(p\). It is naturally a vector space (of dimension \(n\) over \(R\)).
\end{definition}

\printindex

\end{document}